\documentclass[10pt,a4paper]{amsart}
\usepackage[margin=19mm]{geometry}
\usepackage{amsmath,amssymb,mathtools}
\usepackage[scr=rsfs,cal=euler]{mathalfa}
\usepackage{xfrac}
\usepackage[unicode,hidelinks,bookmarks=false]{hyperref}
\newcommand{\ie}{i.e.\ }
\newcommand{\cf}{cf.\ }
\newcommand{\resp}{respectively\ }
\newcommand{\Subsection}{\subsection}
\providecommand{\nobreakdash}{\nobreak\hbox{-}}
\setlength{\emergencystretch}{2em}
\multlinegap0pt
\usepackage{bm}
\usepackage{bbm}
\usepackage{dsfont}
\usepackage{xspace}
\usepackage{cancel}
\usepackage{mathtools}
\usepackage{amsthm}
\makeatletter
\@ifundefined{corollary}{\newtheorem{corollary}{Corollary}}{}
\@ifundefined{theorem}{\newtheorem{theorem}{Theorem}}{}
\@ifundefined{proposition}{\newtheorem{proposition}[theorem]{Proposition}}{}
\makeatother
\def\d{\mathrm{d}}
\newcommand\JP[1]{\langle#1\rangle}
\newcommand\norm[1]{\left\|#1\right\|}
\newcommand{\R}{\mathbb{R}}
\title[Strichartz estimate for discrete Schr\"odinger equation on l]{Strichartz estimate for discrete Schr\"odinger equation on layered King's grid}
\author{Zhiqiang Wan, Heng Zhang}
\date{Source: arXiv:2507.20142v2 (CC BY 4.0)}
\begin{document}
\maketitle

\noindent\emph{Source and licence.} Strichartz estimate for discrete Schr\"odinger equation on layered King's grid. Version arXiv:2507.20142v2 (CC BY 4.0), distributed under the Creative Commons Attribution 4.0 International licence, as recorded in the arXiv licence metadata for this exact version and shown on the abstract page https://arxiv.org/abs/2507.20142v2. Full rights notice: licenses/arxiv-2507.20142-CC-BY-4.0.txt.\par\smallskip

\noindent\textit{Editorial scope.} Counted unit: the Theorem whose source label is \texttt{theorem: POU+decay} in the article (source0.tex lines 555--575, source label \texttt{theorem: POU+decay}), delivered together with the article's own complete original proof of it (source0.tex lines 580--651, one \texttt{proof} environment of 72 lines). No mathematics is written, repaired or re-derived: statement and proof are the article's own text line for line. The proof is detached from the statement in the source: the article states the result at lines 555--575 and prints its proof at lines 580--651, and the proof environment's own header names the statement, so the association is the article's own. Required context retained uncounted: def: General oscillatory integral on the torus (lines 380--382); def: Sigmai (lines 420--426); def: Vi (lines 434--441); coro: decomposition of torus (lines 521--533); prop:fast decay (lines 546--552). Every definition, display and label of those lines is retained unchanged. Omitted from this package and not redistributed: 1--379, 383--419, 427--433, 442--520, 534--545, 553--554, 576--579, 652--1277. Nothing inside a retained excerpt is omitted or altered: every retained body is delivered exactly as the article's lines are written, so no omission receipt of any kind accompanies this package. Citations: the retained excerpts carry no citation command at all, so no bibliography entry of the article is transcribed and no reference is redistributed. Reference adaptation: none was needed and none was made; every reference inside the delivered unit resolves inside it, so no unresolved reference or citation is produced. Printed slips kept literal: no printed word, symbol, hypothesis or number of the counted statement or of its proof was altered. Layout and class: the article's own class is \texttt{amsart} with packages including amssymb, latexsym, amsmath, amsthm, enumitem, mathrsfs, geometry, fullpage, fancyhdr, xcolor, bbm, stmaryrd, hyperref, lineno; the wrapper is the lane's pinned \texttt{amsart} editorial wrapper at 10pt on A4, which restates the theorem-like environments the retained text needs and adds the article's own macro definitions that the retained text needs (\texttt{\textbackslash JP}, \texttt{\textbackslash R}, \texttt{\textbackslash d}, \texttt{\textbackslash norm}), with extra wrapper packages (bm, bbm, dsfont, xspace, cancel, mathtools). The delivered numbering is the unit's own. Assets: no retained span contains a figure environment, a graphics-inclusion command, a PDF-page inclusion, a table or a TikZ picture; no image file is copied into this package, the package declares no asset dependency and no image file is redistributed with it. Licence: version arXiv:2507.20142v2 of this article is distributed under the Creative Commons Attribution 4.0 International licence, as recorded in the arXiv licence metadata for this exact version and shown on the abstract page https://arxiv.org/abs/2507.20142v2. A full rights notice is delivered at licenses/arxiv-2507.20142-CC-BY-4.0.txt. Characters: every retained excerpt is pure ASCII, so the delivered engine, XeLaTeX with its default Latin Modern text font, reports no missing character.
\begin{equation}\label{def: General oscillatory integral on the torus}
    I\left(\xi,\eta;t;\chi\right)=\int_{\left[0,2\pi\right]^{2}}e^{-it\Omega\left(x,y;\xi,\eta\right)}\chi\left(x,y\right)\d x\d y
\end{equation}

\begin{equation}\label{def: Sigmai}
	\begin{aligned}
		\Sigma_{1}&=\left\{\left(x_0,y_0\right)\in\left[0,2\pi\right]^{2}:\det\left(\mathrm{Hess}_{\left(x_0,y_0\right)}\omega\right)\neq0\right\},\\
		\Sigma_{2}&=\left\{\left(x_0,y_0\right)\in\left[0,2\pi\right]^{2}:\left(x_0,y_0\right)\;\mathrm{is\;an}\;A_{2}\;\mathrm{type\;singularity\;of}\;\omega\left(x,y\right)-\left(x,y\right)\cdot\nabla_{\left(x_0,y_0\right)}\omega\right\},\\
		\Sigma_{3}&=\left\{\left(x_0,y_0\right)\in\left[0,2\pi\right]^{2}:\left(x_0,y_0\right)\;\mathrm{is\;an}\;A_{3}\;\mathrm{type\;singularity\;of}\;\omega\left(x,y\right)-\left(x,y\right)\cdot\nabla_{\left(x_0,y_0\right)}\omega\right\}.\\
	\end{aligned}
\end{equation}

\begin{equation}\label{def: Vi}
	\begin{aligned}
        V_{0}&=\left\{v\in\R^{2}:\mathrm{for\;all}\;\left(x_{0},y_{0}\right)\in\left[0,2\pi\right]^{2},\;v\neq\nabla_{\left(x_{0},y_{0}\right)}\omega\right\},\\
		V_{3}&=\left\{v\in\R^{2}:\mathrm{there\;exists\;some}\;\left(x_{0},y_{0}\right)\in \Sigma_{3},\;\mathrm{such\;that}\;v=\nabla_{\left(x_{0},y_{0}\right)}\omega\right\},\\
        V_{2}&=\left\{v\in\R^{2}\setminus V_{3}:\mathrm{there\;exists\;some}\;\left(x_{0},y_{0}\right)\in \Sigma_{2},\;\mathrm{such\;that}\;v=\nabla_{\left(x_{0},y_{0}\right)}\omega\right\},\\
        V_{1}&=\left\{v\in\R^{2}\setminus \left(V_{2}\cup V_{3}\right):\mathrm{there\;exists\;some}\;\left(x_{0},y_{0}\right)\in \Sigma_{1},\;\mathrm{such\;that}\;v=\nabla_{\left(x_{0},y_{0}\right)}\omega\right\}.\\
	\end{aligned}
\end{equation}

\begin{corollary}\label{coro: decomposition of torus}
    Let sets $\{\Sigma_{i}\}_{i=1}^{3}$ defined by \eqref{def: Sigmai} and fix $\left(x_{0},y_{0}\right)\in\left[0,2\pi\right]^{2}$. Then there exist a neighborhood of $\left(x_{0},y_{0}\right)$, denoted as $U_{\left(x_{0},y_{0}\right)}$, such that for all smooth test function $\chi$ supported in $U_{\left(x_{0},y_{0}\right)}$,
    \begin{equation}
        \left|\widetilde{I}\left(\xi,\eta;t;\chi\right)\right|\leq C_{\left(x_{0},y_{0}\right)}\norm{\mathbf{a}\chi}_{C^{3}\left(\R^{2}\right)}\left(1+\left|t\right|\right)^{-3/4}\;\;\mathrm{if}\;\left(x_{0},y_{0}\right)\in \Sigma_{3},
    \end{equation}
    \begin{equation}
        \left|\widetilde{I}\left(\xi,\eta;t;\chi\right)\right|\leq C_{\left(x_{0},y_{0}\right)}\norm{\mathbf{a}\chi}_{C^{3}\left(\R^{2}\right)}\left(1+\left|t\right|\right)^{-5/6}\;\;\mathrm{if}\;\left(x_{0},y_{0}\right)\in \Sigma_{2},
    \end{equation}
    \begin{equation}
        \left|\widetilde{I}\left(\xi,\eta;t;\chi\right)\right|\leq C_{\left(x_{0},y_{0}\right)}\norm{\mathbf{a}\chi}_{C^{3}\left(\R^{2}\right)}\left(1+\left|t\right|\right)^{-1}\;\;\mathrm{if}\;\left(x_{0},y_{0}\right)\in \Sigma_{1},
    \end{equation}
    for all $\left(\xi,\eta\right)\in\R^{2}$.
\end{corollary}

\begin{proposition}\label{prop:fast decay}
     If $\mathrm{supp}\;\chi$ does not contain any critical points of the phase function $\Omega\left(x,y;\xi,\eta\right)$, then for any $M>0$,
    \begin{equation*}
        \left|\widetilde{I}\left(\xi,\eta;t;\chi\right)\right|\leq C\left(M,\chi,d\right)\JP{t}^{-M},
    \end{equation*}
    where $d$ is the infimum of $\left|t^{-1}\left(\xi,\eta\right)-\nabla_{\left(x,y\right)}\omega\right|$ over the support of $\chi$ with respect to $\left(x,y\right)$.
\end{proposition}

\begin{theorem}\label{theorem: POU+decay}
    Let the sets $\{V_{i}\}_{i=0}^{3}$ be defined by \eqref{def: Vi}, $\chi$ be a smooth periodic function on $\left[0,2\pi\right]^{2}$, and the integral $I\left(\xi,\eta;t;\chi\right)$ defined by \eqref{def: General oscillatory integral on the torus}. Then for any fixed $\delta>0$, there exist constant $C_{0}$, $C_{1}$, $C_{2}$ and $C_{3}$ depending on $\chi$ such that
    \begin{itemize}
        \item For all $\left(\xi,\eta\right)$ with $\mathrm{dist}\left(\left(\xi,\eta\right),tV_{3}\right)\leq t\delta$, we have
            \begin{equation}\label{eq:A_{3} decay}
                \left|I\left(\xi,\eta;t;\chi\right)\right|\leq\frac{C_{3}\left(\chi\right)}{\left|t\right|^{3/4}}.
            \end{equation}
        \item For all $\left(\xi,\eta\right)$ with $\mathrm{dist}\left(x,tV_{3}\right)>t\delta$ and $\mathrm{dist}\left(\left(\xi,\eta\right),tV_{2}\right)\leq t\delta$, we have
            \begin{equation}
                \left|I\left(\xi,\eta;t;\chi\right)\right|\leq\frac{C_{2}\left(\chi,\delta\right)}{\left|t\right|^{5/6}}.
            \end{equation}
        \item For all $\left(\xi,\eta\right)$ with $\mathrm{dist}\left(\left(\xi,\eta\right),t\left(V_{2}\cup V_{3}\right)\right)> t\delta$ and $\mathrm{dist}\left(x,tV_{1}\right)\leq t\delta$, we have
            \begin{equation}
                \left|I\left(\xi,\eta;t;\chi\right)\right|\leq\frac{C_{1}\left(\chi,\delta\right)}{\left|t\right|}.
            \end{equation}
        \item For all $\left(\xi,\eta\right)$ with $\mathrm{dist}\left(\left(\xi,\eta\right),t\left(V_{1}\cup V_{2}\cup V_{3}\right)\right)>t\delta$, we have
            \begin{equation}
                \left|I\left(\xi,\eta;t;\chi\right)\right|\leq\frac{C_{0}\left(\chi,\delta,M\right)}{\left|t\right|^{M}}.
            \end{equation}
    \end{itemize}
\end{theorem}

\begin{proof}[Proof of Theorem \ref{theorem: POU+decay}]
    Let $\delta>0$ be a fixed number. When $\left(\xi,\eta\right)\in\R^{2}$ satisfies that $\mathrm{dist}\left(\left(\xi,\eta\right),t\left(V_{1}\cup V_{2}\cup V_{3}\right)\right)>t\delta$, i.e., $\mathrm{dist}\left(t^{-1}\left(\xi,\eta\right),V_{1}\cup V_{2}\cup V_{3}\right)>\delta$, then by taking use of Proposition \ref{prop:fast decay}, we obtain
    \begin{equation*}
        \left|\widetilde{I}\left(\xi,\eta;t;\chi\right)\right|\leq C\left(M,\chi,d\right)\JP{t}^{-M}
    \end{equation*}
    with $d=\delta$. Indeed, we note that $\nabla\omega$ is periodic on $\R^{2}$ and then
    \begin{equation*}
        \begin{aligned}
            &\inf_{\left(x,y\right)\in\mathrm{supp}\left(\mathbf{a}\cdot\chi\right)}\left|t^{-1}\left(\xi,\eta\right)-\nabla_{\left(x,y\right)}\omega\right|=\inf_{\left(x,y\right)\in\mathrm{supp}\left(\chi\right)}\left|t^{-1}\left(\xi,\eta\right)-\nabla_{\left(x,y\right)}\omega\right|>\delta.
        \end{aligned}
    \end{equation*}
    Below we always assume that $\mathrm{dist}\left(\left(\xi,\eta\right),t\left(V_{1}\cup V_{2}\cup V_{3}\right)\right)\leq t\delta$. 
    
    For any $\left(x,y\right)\in\left[0,2\pi\right]^{2}$, by fixing a neighborhood (in particular, we set this neighborhood be a ball) of $\left(x,y\right)$ which satisfies the inequalities in Corollary \ref{coro: decomposition of torus}, we can get a open cover $\{U_{\left(x,y\right)}\}_{\left(x,y\right)\in\left[0,2\pi\right]^{2}}$ of the compact set $\left[0,2\pi\right]^{2}$. Then one may extract a finite sub-cover $\{U_n\}_{1\le n\le N_1}$ and construct a subordinate partition of unity $\{\chi_n\}_{1\le n\le N_1}$ such that $\operatorname{supp}\chi_n\subset U_n$ for each $n$ and every $\chi_n$ satisfies the required periodicity condition. Since
    \begin{equation}
        \left|I\left(\xi,\eta;t;\chi\right)\right|\leq\sum_{1\leq n\leq N_{1}}\left|I\left(\xi,\eta;t;\chi\cdot\chi_{n}\right)\right|=\sum_{1\leq n\leq N_{1}}\left|\widetilde{I}\left(\xi,\eta;t;\chi\cdot\chi_{n}\right)\right|,
    \end{equation}
    and by Corollary \ref{coro: decomposition of torus}, each $\widetilde{I}\left(\xi,\eta;t;\chi\cdot\chi_{n}\right)$ satisfies that (here we use the lowest decay rate $-3/4$)
    \begin{equation*}
        \left|\widetilde{I}\left(\xi,\eta;t;\chi\cdot\chi_{n}\right)\right|\leq C_{n}\norm{\mathbf{a}\chi\cdot\chi_{n}}_{C^{3}\left(\R^{2}\right)}\left(1+\left|t\right|\right)^{-3/4},
    \end{equation*}
    we conclude that
    \begin{equation}\label{eq: -3/4 bound}
        \left|I\left(\xi,\eta;t;\chi\right)\right|\leq C\left(\chi\right)\JP{t}^{-3/4}.
    \end{equation}
    Inequality \eqref{eq: -3/4 bound} holds for any $\left(\xi,\eta\right)\in\R^{2}$, and we will use more detailed assumptions of the open cover to improve this inequality for some particular $\left(\xi,\eta\right)$. 
    
    We define the function $\mathcal{V}:\left[0,2\pi\right]^{2}\to \R^{2}$, where $\mathcal{V}\left(x,y\right)=\nabla_{\left(x,y\right)}\omega$ and hence it's uniformly continuous on $\left[0,2\pi\right]^{2}$. For given $\delta>0$, we choose some small enough $\epsilon>0$ depends on $\delta$ such that
    \begin{equation*}
        \mathrm{diam}\left(\mathcal{V}\left(B_{\epsilon}\right)\right)<\frac{\delta}{2}
    \end{equation*}
    for all ball $B_{\epsilon}$ of radius $\epsilon$. Applying Corollary \ref{coro: decomposition of torus} we get a refine open covering \begin{equation*}
        \left\{U_{\left(x,y\right)}\cap B_{\epsilon}\left(\left(x,y\right)\right)\right\}_{\left(x,y\right)\in\left[0,2\pi\right]^{2}}
    \end{equation*} 
    of $\left[0,2\pi\right]^{2}$, and after choosing a suitable finite sub-cover $\left\{U_{n}\right\}_{1\leq n\leq N_{2}}$, there is a partition of unity $\{\chi_{n}\}_{1\leq n\leq N_{2}}$ such that $\chi_{n}$ supports in $U_{n}$ for each $1\leq n\leq N_{2}$ and each $\chi_{n}$ satisfies the periodic condition. Noting that each subset $U_{n}$, $1\leq n\leq N_{2}$, is identified with a center $\left(x_{n},y_{n}\right)\in\left[0,2\pi\right]^{2}$, we give three index sets to represent those centers in $\Sigma_{i}$,
    \begin{equation*}
        J_{i}:=\left\{1\leq n\leq N_{2}:\left(x_{n},y_{n}\right)\in\Sigma_{i}\right\}\;\mathrm{where}\;i=1,2,3.
    \end{equation*}
    For every $\left(\xi,\eta\right)\in \R^{2}$ with $\mathrm{dist}\left(\left(\xi,\eta\right),tV_{3}\right)>t\delta$, we have
    \begin{equation*}
        \left|t^{-1}\left(\xi,\eta\right)-\nabla_{\left(x,y\right)}\omega\right|>\delta
    \end{equation*}
    for all $\left(x,y\right)\in \Sigma_{3}$ and this implies that $\left|t^{-1}\left(\xi,\eta\right)-\nabla_{\left(x,y\right)}\omega\right|>\delta/2$ for all $\left(x,y\right)\in\cup_{n\in J_{3}}U_{n}$. Proposition \ref{prop:fast decay} gives that for each $n\in J_{3}$, 
    \begin{equation*}
        \left|\widetilde{I}\left(\xi,\eta;t;\chi\cdot\chi_{n}\right)\right|\leq C\left(M,\chi,d\right)\JP{t}^{-M},
    \end{equation*}
    where we can set $d=\delta/2$. Therefore, 
    \begin{equation*}
        \begin{aligned}
            \left|I\left(\xi,\eta;t;\chi\right)\right|\leq&\sum_{1\leq n\leq N_{2}}\left|I\left(\xi,\eta;t;\chi\cdot\chi_{n}\right)\right|=\left(\sum_{n\in J_{3}}+\sum_{n\in J_{1}\cup J_{2}}\right)\left|\widetilde{I}\left(\xi,\eta;t;\chi\cdot\chi_{n}\right)\right|\\
            \leq&C\left(M,\chi,\delta\right)\JP{t}^{-M}+C\left(\chi\right)\JP{t}^{-5/6},
        \end{aligned}
    \end{equation*}
    for each $M\in\mathbb{N}$. In particular, for $M=1$, we have $\left|I\left(\xi,\eta;t;\chi\right)\right|\leq C\left(\chi,\delta\right)\JP{t}^{-5/6}$. 

    Finally, for those $\left(\xi,\eta\right)\in \R^{2}$ with $\mathrm{dist}\left(\left(\xi,\eta\right),t\left(V_{2}\cup V_{3}\right)\right)>t\delta$, we conclude
    \begin{equation*}
        \left|t^{-1}\left(\xi,\eta\right)-\nabla_{\left(x,y\right)}\omega\right|>\delta
    \end{equation*}
    for all $\left(x,y\right)\in \Sigma_{2}\cup \Sigma_{3}$ and this implies that $\left|t^{-1}\left(\xi,\eta\right)-\nabla_{\left(x,y\right)}\omega\right|>\delta/2$ for all $\left(x,y\right)\in\cup_{n\in J_{2}\cup J_{3}}U_{n}$. Using Proposition \ref{prop:fast decay} again, we derive that for each $n\in J_{2}\cup J_{3}$, 
    \begin{equation*}
        \left|\widetilde{I}\left(\xi,\eta;t;\chi\cdot\chi_{n}\right)\right|\leq C\left(M,\chi,d\right)\JP{t}^{-M},
    \end{equation*}
    where we can set $d=\delta/2$. Hence, inequality
    \begin{equation*}
        \begin{aligned}
            \left|I\left(\xi,\eta;t;\chi\right)\right|\leq&\sum_{1\leq n\leq N_{2}}\left|I\left(\xi,\eta;t;\chi\cdot\chi_{n}\right)\right|=\left(\sum_{n\in J_{2}\cup J_{3}}+\sum_{n\in J_{1}}\right)\left|\widetilde{I}\left(\xi,\eta;t;\chi\cdot\chi_{n}\right)\right|\\
            \leq&C\left(M^{\prime},\chi,\delta\right)\JP{t}^{-M^{\prime}}+C\left(\chi\right)\JP{t}^{-1}
        \end{aligned}
    \end{equation*}
    holds for each $M^{\prime}\in\mathbb{N}$. Setting $M^{\prime}=2$, we have $\left|I\left(\xi,\eta;t;\chi\right)\right|\leq C\left(\chi,\delta\right)\JP{t}^{-1}$. We complete the proof.
\end{proof}

\end{document}
